\section{The question: useful context, useful actions}
A company memory may accurately store thousands of records yet fail to supply the few facts needed for a defensible next action. For a marketing research recommendation, those facts can span approved audience policy, an account's current situation, an offering's relevant pain and supporting proof. A recent signal alone is insufficient. Conversely, a capable reader cannot recover company-specific evidence that its context omits.

We ask whether domain-conditioned evidence selection improves \emph{grounded action recommendations}, holding the reader fixed. This separates two interventions often bundled under a ``specialized brain'': changing the evidence presented to a model and changing the model's planning procedure. We measure complete evidence support before reading and action correctness afterward, using the same GPT or Claude within every comparison. The task checks the eligibility of a bounded research action; it does not identify the revenue-maximizing marketing strategy.

The contribution is an auditable controlled comparison, not a new memory architecture: (i) real GPT/Claude calls across seven memory treatments; (ii) native GBrain search and a strong structured-file reference, including a same-backend coverage control; and (iii) an exploratory raw-document versus compiled-wiki check on externally authored company material. The release retains raw responses, unsuccessful attempts, source hashes and executable scoring. Both positive effects and simple baselines that erase them matter for architecture decisions.
